\documentclass[11pt]{article}
\usepackage[margin=1in]{geometry}
\usepackage{booktabs,amsmath,amssymb,graphicx,hyperref,xcolor,url}
\usepackage[numbers]{natbib}
\hypersetup{colorlinks=true,linkcolor=blue,citecolor=blue,urlcolor=blue}
% GENERATED by paper/gen_numbers.py from the claims registry.
% Do not hand-edit. Re-run the generator instead.

\newcommand{\NrowsRaw}{856}
\newcommand{\NrowsDropped}{33}
\newcommand{\MinAccBefore}{20}
\newcommand{\Nrows}{823}
\newcommand{\Nckpt}{38}
\newcommand{\Nfam}{6}
\newcommand{\Nsch}{6}
\newcommand{\MinCkpt}{3}
\newcommand{\NpredNums}{18}
\newcommand{\AdvWorst}{-39.5}
\newcommand{\AdvBadRows}{36}
\newcommand{\AdvBadOver}{15}
\newcommand{\AdvCtrlRows}{9}
\newcommand{\AdvCtrlOver}{2}
\newcommand{\AdvMaxB}{1.5}
\newcommand{\CtrlRatio}{2.6}
\newcommand{\CtrlGap}{-3.78}
\newcommand{\CalibRatio}{171}
\newcommand{\CtrlZmin}{2.9}
\newcommand{\CtrlZmax}{5.3}
\newcommand{\BaseQwenSmall}{47.20}
\newcommand{\PubQwenSmall}{46.83}
\newcommand{\PubQwenSmallBase}{47.57}
\newcommand{\NvfpMarginal}{63}
\newcommand{\AnomQuantMinPct}{23}
\newcommand{\AnomQuantMaxPct}{28}
\newcommand{\NRhCardsListed}{439}
\newcommand{\NRhCardsHarvested}{102}
\newcommand{\NRhCardsWithLibVersion}{12}
\newcommand{\NRhCardsWithActOrder}{33}
\newcommand{\NRhCardsWithDamp}{17}
\newcommand{\NRhCardsNoRecipe}{198}
\newcommand{\NRhHarvestedWithLibVersion}{0}
\newcommand{\NRhHarvestedNoRecipe}{39}
\newcommand{\BaseQwenMid}{59.57}
\newcommand{\PubQwenMid}{60.98}
\newcommand{\ProtoSens}{5.70}
\newcommand{\ExcessN}{6}
\newcommand{\IntervalsExclZero}{0}
\newcommand{\BigLossBelow}{20}
\newcommand{\BigLoss}{27}
\newcommand{\CellsBlocked}{4}
\newcommand{\PooledOneSided}{95.7}
\newcommand{\PubOthersVerifiableCards}{1}
\newcommand{\ProspInside}{121}
\newcommand{\ProspN}{133}
\newcommand{\ProspCov}{91.0}
\newcommand{\ProspCILo}{84.8}
\newcommand{\ProspCIHi}{95.3}
\newcommand{\MaeGain}{0.03}
\newcommand{\TargetMean}{-0.40}
\newcommand{\TargetSd}{1.25}
\newcommand{\NoiseNBench}{11}
\newcommand{\NoiseNThird}{8}
\newcommand{\NoiseMinPct}{15}
\newcommand{\GsmNoiseSd}{2.1}
\newcommand{\BandEmpCov}{91.0}
\newcommand{\BandConfCov}{90.5}
\newcommand{\InfShare}{26}
\newcommand{\InfSharePct}{25.9}
\newcommand{\FiniteConfCov}{88.6}
\newcommand{\FiniteEmpCov}{87.8}
\newcommand{\FiniteConfWidth}{3.97}
\newcommand{\FiniteEmpWidth}{5.53}
\newcommand{\FiniteN}{2240}
\newcommand{\PubIntelPaired}{9}
\newcommand{\NCellsInsuff}{6}
\newcommand{\NCellsTotal}{17}
\newcommand{\NCellsRefused}{0}
\newcommand{\PooledRows}{823}
\newcommand{\PairsPerRow}{5}
\newcommand{\PooledPairs}{4115}
\newcommand{\NSchemesInsuff}{1}
\newcommand{\NSchemesRefused}{0}
\newcommand{\ProspSchemes}{4}
\newcommand{\ProspAllRows}{188}
\newcommand{\ProspAllCov}{91.0}
\newcommand{\ProspClusLo}{85.9}
\newcommand{\ProspClusHi}{94.8}
\newcommand{\ProspClusN}{19}
\newcommand{\ProspRowsWfour}{82}
\newcommand{\ProspCkptsWfour}{11}
\newcommand{\ProspCovWfour}{89.0}
\newcommand{\ProspRowsWfourShare}{62}
\newcommand{\ProspInFamRows}{77}
\newcommand{\ProspInFamInside}{70}
\newcommand{\ProspInFamCov}{90.9}
\newcommand{\ProspInFamCkpts}{7}
\newcommand{\ProspOutFamRows}{56}
\newcommand{\ProspOutFamInside}{51}
\newcommand{\ProspOutFamCov}{91.1}
\newcommand{\ProspOutFamCkpts}{12}
\newcommand{\ProspGemThreeRows}{45}
\newcommand{\ProspGemThreeInside}{40}
\newcommand{\ProspGemThreeCov}{88.9}
\newcommand{\ProspGemThreeCkpts}{8}
\newcommand{\ProspGemThreeGap}{1.1}
\newcommand{\ProspOutSmallRows}{11}
\newcommand{\ProspOutSmallCkpts}{4}
\newcommand{\ProspOutSmallPullup}{2.2}
\newcommand{\ProspCrLo}{86.6}
\newcommand{\ProspCrHi}{95.4}
\newcommand{\ProspCrSe}{2.23}
\newcommand{\ProspNaiveSe}{2.48}
\newcommand{\ProspSeRatio}{0.899}
\newcommand{\ProspIcc}{-0.0061}
\newcommand{\ProspDeff}{1.000}
\newcommand{\ProspRowsPerCkptMax}{14}
\newcommand{\ProspCkptsWithMiss}{10}
\newcommand{\ProspCkptsClean}{9}
\newcommand{\ProspWorstCkptMisses}{3}
\newcommand{\ProspTotalMisses}{12}
\newcommand{\ProspWorstBenchMisses}{5}
\newcommand{\ProspBenchWithMiss}{7}
\newcommand{\ProspWidthCp}{10.48}
\newcommand{\ProspWidthWald}{9.74}
\newcommand{\ProspWidthCr}{8.75}
\newcommand{\ProspNarrowEst}{0.74}
\newcommand{\ProspNarrowUnit}{0.99}
\newcommand{\ProspDiffLo}{-9}
\newcommand{\ProspDiffHi}{+11}
\newcommand{\LThreeSubOneGem}{87.6}
\newcommand{\LThreeSubOneGra}{95.2}
\newcommand{\LThreeSubOneMis}{96.2}
\newcommand{\LThreeSubOneQtf}{89.2}
\newcommand{\LThreeSubOneQth}{100.0}
\newcommand{\LThreeSubTwoGem}{87.5}
\newcommand{\LThreeSubTwoGra}{95.0}
\newcommand{\LThreeSubTwoMis}{100.0}
\newcommand{\LThreeSubTwoQtf}{85.0}
\newcommand{\LThreeSubTwoQth}{100.0}
\newcommand{\LThreeSubThreeGem}{75.5}
\newcommand{\LThreeSubThreeGra}{85.7}
\newcommand{\LThreeSubThreeMis}{91.8}
\newcommand{\LThreeSubThreeQtf}{83.7}
\newcommand{\LThreeSubThreeQth}{93.9}
\newcommand{\LThreeMergedGem}{85.5}
\newcommand{\LThreeMergedGra}{93.5}
\newcommand{\LThreeMergedMis}{96.0}
\newcommand{\LThreeMergedQtf}{87.6}
\newcommand{\LThreeMergedQth}{98.9}
\newcommand{\LofoFallbackPairs}{1091}
\newcommand{\LofoTotalPairs}{4115}
\newcommand{\LofoFallbackRate}{26.5}
\newcommand{\LofoFallRateFp}{60}
\newcommand{\LofoFallRateFpdyn}{20}
\newcommand{\LofoFallRateNvfp}{60}
\newcommand{\LofoFallRateWfour}{0}
\newcommand{\LofoFallRateWsixteen}{59}
\newcommand{\LofoFallRateWint}{20}
\newcommand{\LofoHeldoutPairs}{3024}
\newcommand{\LofoHeldoutCov}{92.0}
\newcommand{\LofoHeldoutOneSided}{96.0}
\newcommand{\PooledCellCov}{91.8}
\newcommand{\ProspRegisteredRepos}{34}
\newcommand{\ProspNameGatedRepos}{9}
\newcommand{\ProspWithGatedInside}{143}
\newcommand{\ProspWithGatedTotal}{156}
\newcommand{\ProspWithGatedCov}{91.7}
\newcommand{\GpqaOtherRows}{0}
\newcommand{\MaxBenchmarksPerRun}{14}
\newcommand{\LofoSchemeMeanWinsFolds}{5}
\newcommand{\LofoSchemeMeanLosesFolds}{1}
\newcommand{\LofoSchemeMeanMin}{0.39}
\newcommand{\LofoSchemeMeanMax}{1.44}
\newcommand{\LofoRidgeMin}{0.40}
\newcommand{\LofoRidgeMax}{1.49}
\newcommand{\LThreeFoldRows}{275}
\newcommand{\LThreeFoldSharePct}{33.4}
\newcommand{\LThreeHoldoutTrainPct}{66.6}
\newcommand{\NonLThreeHoldoutMin}{80}
\newcommand{\NonLThreeHoldoutMax}{93}
\newcommand{\GemmaTwoFoldRows}{59}
\newcommand{\QwenThreeFoldRows}{103}
\newcommand{\NonLThreeFoldRowsMin}{59}
\newcommand{\NonLThreeFoldRowsMax}{162}
\newcommand{\NonLThreeFoldShareMin}{7.2}
\newcommand{\NonLThreeFoldShareMax}{19.7}
\newcommand{\ProspRowsWeightint}{36}
\newcommand{\ProspCkptsWeightint}{5}
\newcommand{\ProspCovWeightint}{91.7}
\newcommand{\ProspRowsFpdyn}{11}
\newcommand{\ProspCkptsFpdyn}{2}
\newcommand{\ProspCovFpdyn}{100.0}
\newcommand{\ProspRowsWeightsixteen}{4}
\newcommand{\ProspCkptsWeightsixteen}{1}
\newcommand{\ProspCovWeightsixteen}{100.0}
\newcommand{\MaeFloor}{0.507}
\newcommand{\HeadroomPp}{0.262}
\newcommand{\PredBestHeadroomPct}{17.8}
\newcommand{\SchemeHeadroomPct}{13.0}
\newcommand{\NrowsGpqaMain}{6}
\newcommand{\NrowsGpqaMainNorm}{6}
\newcommand{\NrowsGpqaMainCot}{2}
\newcommand{\NrowsGpqaDiamond}{8}
\newcommand{\NrowsGpqaDiamondCot}{3}
\newcommand{\GbSensNitems}{0.0093}
\newcommand{\GbSensNitemsRows}{6}
\newcommand{\GbSensMmluPct}{23}
\newcommand{\GbSensMmluRows}{2}
\newcommand{\MaeRidgeTuned}{0.7248}
\newcommand{\MaeGradboostTuned}{0.7228}
\newcommand{\TuneShiftRidge}{0.0707}
\newcommand{\TuneShiftGb}{0.0628}
\newcommand{\TuneRidgeAlphaMax}{1000}
\newcommand{\TuneRidgeFoldsEdge}{5}
\newcommand{\TuneGbFoldsMinLr}{6}
\newcommand{\TuneGbFoldsMaxLtwo}{6}
\newcommand{\TuneOuterFolds}{6}
\newcommand{\TuneGbGrid}{81}
\newcommand{\TuneRidgeGrid}{11}
\newcommand{\TuneSchemeVsRidge}{+0.0106}
\newcommand{\TuneSchemeVsRidgeLo}{-0.0049}
\newcommand{\TuneSchemeVsRidgeHi}{+0.0247}
\newcommand{\TuneSchemeVsGb}{+0.0126}
\newcommand{\TuneSchemeVsGbLo}{-0.0117}
\newcommand{\TuneSchemeVsGbHi}{+0.0388}
\newcommand{\TuneGlobalVsRidgeLo}{+0.0252}
\newcommand{\TuneGlobalVsRidgeHi}{+0.0648}
\newcommand{\MaeGainLo}{0.01}
\newcommand{\MaeGainHi}{0.06}
\newcommand{\MaeGainFamsPos}{5}
\newcommand{\NoiseMaxPct}{119}
\newcommand{\NoiseNOver}{1}
\newcommand{\NoiseNSmall}{5}
\newcommand{\RefuseBelow}{85}
\newcommand{\BootLoWsixteenBig}{47}
\newcommand{\BootHiWsixteenBig}{100}
\newcommand{\BootLoWfourSmall}{68.0}
\newcommand{\BootHiWfourSmall}{83.3}
\newcommand{\PubChecked}{10}
\newcommand{\PubCards}{260}
\newcommand{\PubOthersCards}{222}
\newcommand{\PubOthersCount}{9}
\newcommand{\PubOthersPaired}{10}
\newcommand{\PubOthersVerified}{6}
\newcommand{\PubCtrlVerified}{147}
\newcommand{\LoFpeight}{-1.11}
\newcommand{\HiFpeight}{+0.48}
\newcommand{\WorstFpeight}{-1.90}
\newcommand{\SevFpeight}{0.0}
\newcommand{\NFpeight}{80}
\newcommand{\LoWeightint}{-1.81}
\newcommand{\HiWeightint}{+1.13}
\newcommand{\WorstWeightint}{-4.04}
\newcommand{\SevWeightint}{1.5}
\newcommand{\NWeightint}{195}
\newcommand{\LoFpdyn}{-1.44}
\newcommand{\HiFpdyn}{+1.34}
\newcommand{\WorstFpdyn}{-8.72}
\newcommand{\SevFpdyn}{1.0}
\newcommand{\NFpdyn}{200}
\newcommand{\LoWeightsixteen}{-0.98}
\newcommand{\HiWeightsixteen}{+0.81}
\newcommand{\WorstWeightsixteen}{-3.04}
\newcommand{\SevWeightsixteen}{1.2}
\newcommand{\NWeightsixteen}{86}
\newcommand{\LoWfour}{-2.87}
\newcommand{\HiWfour}{+1.40}
\newcommand{\WorstWfour}{-8.86}
\newcommand{\SevWfour}{5.7}
\newcommand{\NWfour}{194}
\newcommand{\LoNvfp}{-4.18}
\newcommand{\HiNvfp}{+1.87}
\newcommand{\WorstNvfp}{-8.12}
\newcommand{\SevNvfp}{14.7}
\newcommand{\NNvfp}{68}
\newcommand{\MaeSchememean}{0.7354}
\newcommand{\MaeGlobalmean}{0.7694}
\newcommand{\MaeSchemebench}{0.769}
\newcommand{\MaeRidge}{0.796}
\newcommand{\MaeBenchmean}{0.778}
\newcommand{\MaeGradboost}{0.786}
\newcommand{\CovCellWfourSmall}{76.4}
\newcommand{\OneSidedCellWfourSmall}{76.4}
\newcommand{\CkptsCellWfourSmall}{2}
\newcommand{\RowsCellWfourSmall}{11}
\newcommand{\CovCellWsixteenBig}{73.5}
\newcommand{\OneSidedCellWsixteenBig}{81.3}
\newcommand{\CkptsCellWsixteenBig}{5}
\newcommand{\RowsCellWsixteenBig}{31}


\title{\textbf{Calibrated Risk Envelopes for LLM Quantization,\\
and the Limits of Predicting Compression Damage}}
\author{Grace Jackson}
\date{September 2026}

\begin{document}
\maketitle

\begin{abstract}
\noindent
Quantizing a large language model reduces serving cost and costs some accuracy.
Determining how much normally requires re-running a full benchmark suite, which
is the expense quantization was meant to avoid. We ask whether the accuracy
delta can be predicted instead. Using \Nrows{} published evaluations spanning
\Nckpt{} base checkpoints, \Nfam{} model families and \Nsch{} quantization
schemes, we report three results. First, a \emph{negative} one: per-model
prediction carries almost no signal. Under leave-one-family-out validation the
best predictor, a per-scheme mean, beats a global mean by \(\MaeGain\)pp of MAE (95\% checkpoint-bootstrap CI \([\MaeGainLo, \MaeGainHi]\)), and both ridge
regression and gradient boosting are \emph{worse} than guessing the average,
because evaluation noise accounts for at least a third of the variance on \NoiseNThird{} of \NoiseNBench{} benchmarks. The features we tested are
card-level (bit-width, scheme, size, benchmark identity, headroom); we did
not test predictors that use checkpoint-derived features such as weight
statistics or activation-outlier metrics.
Second, a positive one: conditioning only on the quantization scheme yields
per-scheme conformal-style intervals with \(\ProspCov\%\) two-sided empirical
coverage (\(\ProspInside/\ProspN\), 95\% CI \([\ProspCILo,\ProspCIHi]\)) on
unseen checkpoints. Under the family definition of \S\ref{sec:data} the
strict set splits into \ProspInFamCkpts{} in-family checkpoints
(DeepSeek-R1-Distill of Llama or Qwen bases) at \ProspInFamCov\% and
\ProspOutFamCkpts{} out-of-family at \ProspOutFamCov\%; the
checkpoint-cluster 95\% CI on the difference is [\ProspDiffLo, \ProspDiffHi]pp
and does not resolve whether the two are equivalent. Within out-of-family,
Gemma-3 alone (the only sub-group with more than 10 rows) covers at
\ProspGemThreeCov\% (\ProspGemThreeInside/\ProspGemThreeRows), \ProspGemThreeGap{}pp below
nominal. Those checkpoints cover \ProspSchemes{} of the \Nsch{} schemes,
and W4A16 supplies \ProspRowsWfour{} of the \ProspN{} rows.
Each (scheme, size-band) cell returns one of three verdicts: trusted; insufficient evidence, where the interval is still
shown but the scheme is demoted and flagged; or refused, where the interval is withheld. Currently \NCellsInsuff{} of \NCellsTotal{}
cells are insufficient evidence; on this corpus the third state, as defined
in \S\ref{sec:method}, does not fire. Third, and most consequential for
practice: the published corpus
is outcome-truncated. Every configuration in it is one that someone chose to
release. We rented an A100 and measured deliberately-faulted configurations,
observing losses to \(\AdvWorst{}\)pp against a published worst case of
\(\WorstWfour{}\)pp---measured, we stress, only on models of \AdvMaxB{}B and
below, and on configurations engineered to fail rather than typical ones. This is an existence proof that far worse outcomes
exist than the corpus contains; it does not estimate how common they are. We argue the resulting artifact should be described as a
calibrated historical baseline rather than a predictor, and we state precisely
why its coverage guarantee does not transfer to a practitioner's own untuned
recipe. Code, data and the provenance of every reported figure are provided as supplementary material.
\end{abstract}

\section{Introduction}

Post-training quantization is now routine. A practitioner choosing between
W4A16, W8A8 and FP8 wants one number: how much accuracy will this cost? In
practice the only way to know is to run the evaluation.

This paper asks whether that number can be forecast from the public record.
Model publishers release quantized checkpoints with benchmark tables comparing
the compressed model against its dense parent. Aggregated, these are hundreds
of paired measurements of exactly the quantity of interest (\NrowsRaw{} extracted in our harvest).

Our contributions:

\begin{enumerate}
  \item \textbf{A negative result} (\S\ref{sec:negative}). Per-model
    prediction of the accuracy delta \emph{from card-level features} does
    not work, and we show why: the target is dominated by evaluation noise.
    Checkpoint-derived features (weight statistics, activation-outlier
    metrics) were not tested. We report this as the primary finding rather
    than burying it.
  \item \textbf{A calibrated envelope that does work} (\S\ref{sec:results}),
    conditioned on scheme alone, with measured coverage on held-out checkpoints
    and a three-state verdict for each (scheme, size-band) cell: \emph{trusted};
    \emph{insufficient evidence}, where the interval is still shown but the
    scheme is demoted and the cell flagged as one the tool cannot judge; or
    \emph{refused}, where the interval is withheld. Currently \NCellsInsuff{} of
    \NCellsTotal{} cells are insufficient evidence; on this corpus the third
    state, as defined in \S\ref{sec:method}, does not fire (see \S\ref{sec:method}
    for why).
  \item \textbf{An existence proof that the corpus omits severe outcomes}
    (\S\ref{sec:tail}), via GPU experiments on deliberately-faulted
    configurations---data the published record structurally cannot contain.
    It shows that losses far beyond anything published are reachable on small
    models. It does not estimate how large the bias is: our attempt at a
    correction (in the supplement) did not work well enough to ship.
  \item \textbf{An audit} (\S\ref{sec:audit}) recording errors we found in
    our own work, including ones that changed published numbers and claims.
\end{enumerate}

\section{Related Work}

\paragraph{Predicting evaluations.} BenchPress~\cite{zeng2026benchpress}
establishes the mechanism closest to ours: predict an eval result, estimate
whether to trust it, wrap it in a conformal interval, and decide whether to skip
the real run. Their score matrix over 84 frontier models and 133 benchmarks is
approximately rank-2. We do not claim novelty against them on mechanism. Their
estimand is a model's absolute score; ours is the paired delta between a model
and its compressed self, and quantized checkpoints do not appear in their paper.

\paragraph{Conformal prediction on compressed models.}
\citet{tong2026compression} apply conformal prediction directly to
quantized and sparse LLMs, which is a closer domain than BenchPress. The
estimand differs: they evaluate predictive uncertainty (coverage and set
size) \emph{after} running the compressed model. Ours is constructed over
historical accuracy deltas to give an interval \emph{before} running anything.
Their Figure~2 plots
the same quantity we predict, \(\text{Acc}_{\text{compressed}} -
\text{Acc}_{\text{dense}}\), but measures it. Their finding that compression
decouples accuracy from uncertainty bounds what an accuracy-only estimate like
ours can claim (\S\ref{sec:limits}).

\paragraph{Quantization methods.} The checkpoints we harvest are produced with
GPTQ~\cite{frantar2022gptq} and SmoothQuant~\cite{xiao2022smoothquant} via
llm-compressor~\cite{llmcompressor}. \citet{kumar2024scaling} fit a functional form that predicts degradation
from training and inference in varied precisions, using pretraining runs of models up to 1.7B parameters. We fit no such law: our
envelope is empirical, per scheme, over published post-training-quantized checkpoints.

\paragraph{The source of our data.} \citet{kurtic2024halfmillion}
report over 500{,}000 evaluations of quantized Llama-3.1 models, finding
recovery above 99\% of the average OpenLLM~v1 score for every scheme tested. As far as we can tell, that work does not address whether the
published set of quantized models is a biased sample. Characterising that gap is
the contribution of \S\ref{sec:tail}.

\paragraph{Selection bias.} Survivorship bias has been
identified and measured in a widely used retrieval benchmark~\cite{gupta2022survivorship}.
Conformal prediction's guarantee is known to fail under selection:
conditional on a selection event, calibration data are no longer exchangeable
with a test point~\cite{barber2023beyond,jin2024focal}. We apply this directly
to our own artifact in \S\ref{sec:limits}.

\section{Data}\label{sec:data}

We harvest published model cards into rows of
\((\text{model}, \text{scheme}, \text{benchmark},
\text{acc}_{\text{before}}, \text{acc}_{\text{after}})\), yielding \NrowsRaw{}
rows, of which \Nrows{} remain after dropping \NrowsDropped{} whose baseline accuracy is below
\MinAccBefore{}\% (near chance, where a delta is uninformative), across \Nckpt{} base checkpoints, \Nfam{} families and
\Nsch{} schemes.

\paragraph{An integrity gate.} Cards typically print a \emph{Recovery \%}
column alongside the before/after pair. This is redundant information, and we
exploit it: for every three-column row we check
\(\text{recovery} \approx 100 \times \text{after}/\text{before}\) to a tolerance
derived from the card's own printed precision, and reject rows that disagree.
This caught spurious values that our parser produced under column-order variation between
card templates, and internally inconsistent cards. Rejections are recorded rather
than silently dropped.

\paragraph{A structural limitation, stated at the outset.} All rows come from a
single publisher using a single toolchain. Holding out a model family tests
transfer across \emph{models}; nothing in this design tests transfer across
\emph{quantization practitioners}. \S\ref{sec:tail} gives direct evidence that
this distinction matters.

\paragraph{What we mean by family.} Throughout, two checkpoints share
a family when one is derived from the other's pretraining weights (via
distillation, continued pretraining, quantization, or fine-tuning), or when
they share a common pretraining base checkpoint. A distillation is a member
of its teacher's family; a new pretraining run is a new family, regardless
of shared name. This is the mechanistic definition: quantization damage
depends on the weight distribution and activation-outlier structure the
checkpoint inherits, and derivative products inherit them, while new
pretraining runs produce different weight distributions and different
quantization behaviour. In particular, under this definition Gemma-3 and
Gemma-2 are separate families (independent pretraining runs), Llama-3.1,
Llama-3.2 and Llama-3.3 all belong to the same Llama-3 family
(Llama-3.2 1B/3B are pretraining-time distillations from Llama-3.1
8B/70B logits; Llama-3.2 11B/90B are Llama-3.1 plus a vision adapter;
Llama-3.3-70B is a Llama-3.1-70B fine-tune -- Meta's own model cards
state each of these lineages), and DeepSeek-R1-Distill-Llama-8B is a
member of that Llama-3 family (it uses Llama-3.1-8B weights as
initialization). The \Nfam{} training families are Gemma-2, Granite,
Llama-3, Mistral, Qwen-2.5 and Qwen-3. An earlier version of this paper
counted Llama-3.1, Llama-3.2 and Llama-3.3 as three separate families, so
the corpus was reported as having 8 training families rather than 6, and
leave-one-family-out was run over 8 folds rather than 6; the audit note
in \S\ref{sec:audit} records the change and its effect on the aggregate
LOFO figure.

\paragraph{The pathway is only sparsely documented.} A live scan on
2026-09-26 across the \NRhCardsListed{} RedHatAI quantized-model cards
currently published on Hugging Face (of which \NRhCardsHarvested{}
contribute rows to our harvested dataset) shows how sparsely the pathway
that produced these outcomes is documented: \NRhCardsWithLibVersion{} of
\NRhCardsListed{} cards pin a quantization-library version,
\NRhCardsWithActOrder{} mention activation reordering,
\NRhCardsWithDamp{} mention dampening, and \NRhCardsNoRecipe{} of
\NRhCardsListed{} report none of \texttt{num\_calibration\_samples},
\texttt{max\_sequence\_length}, activation reordering or dampening. On the
\NRhCardsHarvested{}-card subset that our dataset draws rows from,
\NRhHarvestedWithLibVersion{} pin a library version and
\NRhHarvestedNoRecipe{} of \NRhCardsHarvested{} report none of those four
parameters. (This RedHatAI-only scan is distinct from the \PubChecked{}-publisher,
\PubCards{}-card publisher census reported in \S\ref{sec:audit}, which is an
earlier stable snapshot spanning multiple publishers and asks a different
question.) The published record we calibrate against is therefore a record
of outcomes with the pathway that produced them substantially undocumented
on the publisher whose cards contribute the largest single share of our
corpus, which is one of the reasons an outcome-truncation argument
(\S\ref{sec:tail}) cannot be closed from the corpus alone.

\section{Method}\label{sec:method}

\paragraph{Conformal-style intervals.} We use the residual-quantile
construction of split conformal regression~\cite{lei2018distribution}: given
residuals \(\{s_i\}_{i=1}^{n}\) and level \(\alpha\), the interval half-width
is the \(\lceil (n+1)(1-\alpha) \rceil\)-th smallest score, which under a
proper calibration/test split gives exact finite-sample coverage under
exchangeability. Where \(n\) is too small for that index to exist, we return
\(\infty\) rather than a misleadingly finite band. The shipped per-scheme
interval fits its centre and its half-width on the \emph{same} rows (an
in-sample residual-quantile band), so the split-conformal guarantee holds only
for the explicit calibration-set variant used inside the leave-one-family-out
audit; the empirical prospective figure below is what carries the shipped
intervals.

\paragraph{Mondrian (group-conditional) calibration.} Calibrating separately within categories is the Mondrian construction of \citet{vovk2005algorithmic}, applied to regression by \citet{bostrom2020mondrian}. A single marginal
interval advertised at 90\% covers only \NvfpMarginal\% of NVFP4 rows
under leave-one-family-out (recomputed live)---the users who most need it.
We therefore calibrate per scheme.

\paragraph{A three-state verdict.} For each (scheme, size-band) cell we
measure coverage directly, and the tool returns one of three verdicts. A cell
is \emph{trusted} when measured coverage is adequate and adequately supported.
It is \emph{refused}---\texttt{INSUFFICIENT CALIBRATION}, no interval---when
coverage is measurably poor: the checkpoint-bootstrap interval on its coverage
lies wholly below the \RefuseBelow\% line, on at least \MinCkpt{} checkpoints. It is
\emph{insufficient evidence} when the cell rests on
fewer than \MinCkpt{} distinct checkpoints, or when a checkpoint-level cluster
bootstrap 90\% interval on its coverage straddles the \RefuseBelow\% line. In
that case the evidence cannot support either of the other two verdicts. The
tool still prints the all-sizes interval for such a cell, but demotes the scheme
to the lowest tier and states that the cell cannot be judged; only a refused
cell has its interval withheld.

Coverage is scored \emph{one-sided}: how often the true delta stayed at or
above the interval's lower bound. The risk this tool bounds is accuracy loss,
and a model that beats its envelope has exposed no one to anything, so counting
that as a failure would flag cells that are merely outperforming. The
interval is symmetric, so a 90\% two-sided nominal corresponds to a
one-sided nominal near 95\%; pooled one-sided coverage across all
\PooledRows{} rows, with each row scored under \PairsPerRow{} different
calibration families, is \PooledOneSided\%---essentially at that target,
not above it.

Currently \NCellsRefused{} of \NCellsTotal{} cells are refused. Two cells were flagged under the earlier two-state rule, and neither is a
confident refusal under the three-state one. \texttt{w4a16|<2B} had one-sided
coverage of \OneSidedCellWfourSmall\% over \RowsCellWfourSmall{} rows
(two-sided \CovCellWfourSmall\%), but those rows come from only
\CkptsCellWfourSmall{} checkpoints. \texttt{w8a16|>10B} had \OneSidedCellWsixteenBig\%
over \RowsCellWsixteenBig{} rows from \CkptsCellWsixteenBig{} checkpoints
(two-sided \CovCellWsixteenBig\%), and its bootstrap interval,
[\BootLoWsixteenBig, \BootHiWsixteenBig]\%, straddles the threshold. Both are
now \emph{insufficient evidence}. We had scoped the change to those two cells;
applying the rule to every cell moved \NCellsInsuff{} of \NCellsTotal{} into
that state, including several whose point-estimate coverage looked comfortable.
The same measurement and the same rule judge each scheme's all-sizes coverage: \NSchemesInsuff{} of
the \Nsch{} schemes is insufficient evidence at that level and \NSchemesRefused{} is refused.
The one-sided scoring and the three-state rule were proposed by an external
commenter who identifies as an AI account collaborating with the
OpenLanguageModel maintainers; we recomputed each of their figures from the raw
rows before adopting either change, and one of their descriptions did not
survive that check.

\paragraph{Post-hoc thresholds, honestly labeled.} The three-state rule was
proposed for two specific cells and then applied to every cell; the
particular thresholds---\RefuseBelow\% for the refusal line, \MinCkpt{}
checkpoints for the support floor, a 90\% cluster bootstrap for the
straddle test---were not derived from a first-principles argument, they
were the values under which those two cells came out the way we thought
they should. On this corpus one further consequence of that choice is
that the \emph{refused} state, as defined, does not fire: every cell with
enough checkpoints to reach the second gate has a bootstrap upper bound
comfortably above \RefuseBelow\%, and the one cell whose upper bound
does sit below it (\texttt{w4a16|<2B}, upper bound \BootHiWfourSmall\%)
has too few checkpoints to reach that gate and lands in
\emph{insufficient evidence} instead. So the current three states are, on
this corpus, effectively two: \emph{trusted} and \emph{insufficient
evidence}. This is a design finding, not a runtime bug, and we state it
here rather than let a reviewer discover it.

\paragraph{A support floor.} Coverage measured on many benchmarks from one
checkpoint is not independent evidence; it is one quantization run viewed many
ways. Cells resting on fewer than \MinCkpt{} distinct checkpoints cannot reach
the top tier regardless of measured coverage. \CellsBlocked{} cells are affected.

\section{The negative result}\label{sec:negative}

Under leave-one-family-out validation, row-weighted MAE:

\begin{center}
\begin{tabular}{lcc}
\toprule
predictor & MAE (pp) & beats global mean? \\
\midrule
per-scheme mean (shipped) & \MaeSchememean & yes \\
global mean (baseline)    & \MaeGlobalmean & --- \\
per-(scheme \(\times\) benchmark) & \MaeSchemebench & \textbf{no} \\
ridge regression        & \MaeRidge      & \textbf{no} \\
per-benchmark mean        & \MaeBenchmean  & \textbf{no} \\
gradient boosting         & \MaeGradboost  & \textbf{no} \\
\bottomrule
\end{tabular}
\end{center}

The per-scheme mean's advantage is small: \MaeGain{}pp, with a checkpoint-bootstrap 95\% CI of \([\MaeGainLo, \MaeGainHi]\)pp, and it is
positive in \MaeGainFamsPos{} of the \Nfam{} held-out families. It is a distinguishable but modest gain.
Ridge regression (\(\alpha=3.0\), standardised features, no cross-validated
tuning) and histogram gradient boosting (\texttt{max\_depth=3},
\texttt{max\_iter=300}, \texttt{learning\_rate=0.05},
\texttt{min\_samples\_leaf=15}, \texttt{l2\_regularization=1.0},
\texttt{random\_state=0}, no cross-validated tuning) are both \emph{worse}
than predicting the global average. Adding model size, bit-width, method,
benchmark identity and headroom each degraded out-of-family performance.
All features tested are card-level; we did not compute any checkpoint-level
statistic (weight kurtosis, activation outliers, singular-value structure)
and feed it in as a predictor. Those directions are open.

\paragraph{Why.} Estimating per-benchmark evaluation noise from near-lossless
schemes---whose true delta should be approximately zero---gives a noise floor
that accounts for anywhere from \NoiseMinPct\% of the observed variance to all of it across \NoiseNBench{} benchmarks. On GSM8K
the deltas of supposedly lossless schemes have a standard deviation of \GsmNoiseSd{}pp, while the quantity being predicted has mean
\(\TargetMean\)pp. The target is a difference of two noisy quantities that largely
cancel: the numerically smallest and hardest-to-predict summary of quantization
damage available. Two cautions on the noise figures. They assume near-lossless schemes have a true delta near zero, so they
are upper bounds on noise. And \NoiseNSmall{} of the \NoiseNBench{} estimates rest on 10 or fewer rows; one, at \NoiseMaxPct\%, exceeds the total
variance, which we attribute to sampling error rather than treat as a finding.

\paragraph{Implication.} From card-level features, the entire fitted
artifact is \NpredNums{} numbers.
A raw per-scheme historical quantile band, with no model and no conformal
machinery, is a serious competitor: it looked better covered in-sample, though under the strict protocol that advantage disappears (\S\ref{sec:audit}). What conformal contributes on the calibration-set variant used inside the leave-one-family-out audit is a finite-sample guarantee under exchangeability; the shipped intervals fit their centre and half-width on the same rows and do not carry that guarantee (\S\ref{sec:method}). What the machinery contributes to the shipped artifact is a principled rule for small-\(n\) schemes, not accuracy. We therefore describe the result as a \emph{calibrated historical
baseline}, not a predictor.

\section{Results}\label{sec:results}

\paragraph{What this tool is for.} The scenario it answers is: a
practitioner is considering a scheme on a checkpoint from a family already
represented in the published record and wants a calibrated range for the
accuracy delta with an honest verdict on whether the record supports that
range for that cell. Within that scope, each cell returns an interval and
one of three tier labels; outside it, the tool declines to speak. It is
not a scheme-safety certificate, does not bound a scheme's worst case,
and does not extrapolate to recipes or families the corpus has not
already seen -- \S\ref{sec:tail} and \S\ref{sec:limits} state each of
those out-of-scope cases directly.

\begin{center}
\begin{tabular}{lcccc}
\toprule
scheme & 90\% interval (pp) & worst observed & \% losing \(>\)3pp & evals \\
\midrule
fp8         & [\LoFpeight, \HiFpeight]           & \WorstFpeight{}pp      & \SevFpeight\%      & \NFpeight \\
w8a8\_int   & [\LoWeightint, \HiWeightint]       & \WorstWeightint{}pp    & \SevWeightint\%    & \NWeightint \\
fp8\_dynamic& [\LoFpdyn, \HiFpdyn]               & \WorstFpdyn{}pp        & \SevFpdyn\%        & \NFpdyn \\
w8a16       & [\LoWeightsixteen, \HiWeightsixteen]& \WorstWeightsixteen{}pp& \SevWeightsixteen\%& \NWeightsixteen \\
w4a16       & [\LoWfour, \HiWfour]               & \WorstWfour{}pp        & \SevWfour\%        & \NWfour \\
nvfp4       & [\LoNvfp, \HiNvfp]                 & \WorstNvfp{}pp         & \SevNvfp\%         & \NNvfp \\
\bottomrule
\end{tabular}
\end{center}

\paragraph{Prospective validation.} After freezing the method we harvested
new checkpoints not used during development. The strict protocol
(\texttt{verify/independent\_check.py}) excludes a prospective card on
one of two grounds: any card whose family label overlaps a training
family (Llama-3.* as one family under Definition B, and Qwen-3), or
any card whose format informed a parser fix (Llama-4). The two grounds
are separate; only Llama-4 is excluded for the parser reason. Under
this protocol two-sided coverage
is \(\ProspInside/\ProspN = \ProspCov\%\), 95\% Clopper--Pearson CI
\([\ProspCILo, \ProspCIHi]\), against a nominal 90\% two-sided level.
Note the layering: this filter runs by name-pattern regex and does not
implement the mechanistic family definition of \S\ref{sec:data}. Under
that definition, \ProspInFamCkpts{} of the \ProspClusN{} strict
checkpoints (all DeepSeek-R1-Distill of a Llama-3.* or Qwen-2.5 base)
are in-family derivatives that the name-pattern filter let through. The pooled figure therefore mixes ``passed the strict
filter'' with ``is in-family under \S\ref{sec:data}''. The
family-split reporting below applies \S\ref{sec:data} to what survives
the filter and states both sides separately. The \ProspN{} strict rows
come from \ProspClusN{} distinct checkpoints, which we hand-classify
against the training-family list. Under the
definition of \S\ref{sec:data}: \ProspInFamCkpts{} in-family checkpoints
(all DeepSeek-R1-Distill of a Llama-3.* or Qwen-2.5 base in
training) contribute \ProspInFamRows{} rows at \ProspInFamCov\% two-sided
coverage; \ProspOutFamCkpts{} out-of-family checkpoints (Gemma-3, SmolLM
and SmolLM3, NVIDIA-Nemotron-Nano-9B) contribute \ProspOutFamRows{} rows
at \ProspOutFamCov\%. Both point estimates land near the nominal 90\%
two-sided level, but the checkpoint-cluster 95\% CI on the in-minus-out
difference is [\ProspDiffLo, \ProspDiffHi]pp: at this sample size we
cannot resolve whether in-family and out-of-family coverage are
equivalent, and the paper does not claim they are. The aggregate
out-of-family figure carries the same Simpson caveat as the pooled
headline: Gemma-3 alone (the only sub-group with more than 10 rows)
covers at \ProspGemThreeCov\% (\ProspGemThreeInside/\ProspGemThreeRows
rows from \ProspGemThreeCkpts{} checkpoints), \ProspGemThreeGap{}pp below nominal, and
SmolLM/SmolLM3 and Nemotron-Nano contribute \ProspOutSmallRows{} rows at
100\% from \ProspOutSmallCkpts{} small-sample checkpoints, which pull the
out-of-family aggregate up by \ProspOutSmallPullup{}pp above Gemma-3
alone. The classification is enumerated for these \ProspClusN{}
checkpoints, not a general algorithm (\S\ref{sec:audit}). A from-scratch, standard-library-only reimplementation sharing no
code with the pipeline reproduces this figure with zero disagreements.
That checks the computation, not the validity of the design.

\paragraph{Merged-fold disclosure: Llama-3.1 vs 3.2 vs 3.3.} Under
Definition B the six-fold LOFO merges Llama-3.1, Llama-3.2 and Llama-3.3
into one Llama-3 fold. The merged mean can conceal that Llama-3.3
alone consistently sits several points below its cousins under every
calibration family, so we break out the held-out Llama-3 fold by
subgroup. Rows are the calibration family; columns are the three
subgroups the merged fold contains.
\begin{center}
\begin{tabular}{lrrrr}
\toprule
calibration & Llama-3.1 & Llama-3.2 & Llama-3.3 & merged \\
\midrule
gemma-2  & \LThreeSubOneGem\% & \LThreeSubTwoGem\% & \LThreeSubThreeGem\% & \LThreeMergedGem\% \\
granite  & \LThreeSubOneGra\% & \LThreeSubTwoGra\% & \LThreeSubThreeGra\% & \LThreeMergedGra\% \\
mistral  & \LThreeSubOneMis\% & \LThreeSubTwoMis\% & \LThreeSubThreeMis\% & \LThreeMergedMis\% \\
qwen2.5  & \LThreeSubOneQtf\% & \LThreeSubTwoQtf\% & \LThreeSubThreeQtf\% & \LThreeMergedQtf\% \\
qwen3    & \LThreeSubOneQth\% & \LThreeSubTwoQth\% & \LThreeSubThreeQth\% & \LThreeMergedQth\% \\
\bottomrule
\end{tabular}
\end{center}
Under every calibration family Llama-3.3 sits below both Llama-3.1 and
Llama-3.2, and the merged mean is pulled up by the two better-covered
subgroups. Under an 8-family reporting this weakness would show as
Llama-3.3's own held-out fold at 85.1\% (the worst held-out family).
Under Definition B it is folded into 92.2\% on the merged Llama-3 mean;
we report the breakdown here so the merge does not hide it.

\paragraph{What the prospective set does and does not cover.} The \ProspN{} strict rows come from \ProspClusN{} quantized checkpoints,
not \ProspN{} independent observations: a bootstrap that resamples checkpoints gives a 95\% interval of
\([\ProspClusLo, \ProspClusHi]\)\%, close to the row-level one. The rows cover only \ProspSchemes{} of the \Nsch{} schemes; FP8 and NVFP4
have no prospective rows, so the figure says nothing about them.

\begin{center}
\begin{tabular}{lccc}
\toprule
scheme & strict rows & checkpoints & covered (\%) \\
\midrule
w4a16       & \ProspRowsWfour & \ProspCkptsWfour & \ProspCovWfour \\
w8a8\_int   & \ProspRowsWeightint & \ProspCkptsWeightint & \ProspCovWeightint \\
fp8\_dynamic& \ProspRowsFpdyn & \ProspCkptsFpdyn & \ProspCovFpdyn \\
w8a16       & \ProspRowsWeightsixteen & \ProspCkptsWeightsixteen & \ProspCovWeightsixteen \\
\bottomrule
\end{tabular}
\end{center}

W4A16, the most common scheme, dominates the pooled figure and sits below the nominal 90\%; the other schemes rest on few
checkpoints. Over all \ProspAllRows{} harvested prospective rows, before the strict filter, coverage is \ProspAllCov\%, so the filter
did not manufacture the headline. That said, the pooled \ProspCov\% is
pulled up by \ProspRowsFpdyn{} FP8-dynamic rows from \ProspCkptsFpdyn{}
checkpoints and \ProspRowsWeightsixteen{} W8A16 rows from \ProspCkptsWeightsixteen{}
checkpoint, both at \ProspCovFpdyn\% coverage: W4A16, at \ProspRowsWfour{}
of the \ProspN{} strict rows (\ProspRowsWfourShare\%), comes in at
\ProspCovWfour\%, which is below the nominal 90\%. The pooled number
should be read as an average across schemes with very different sample
sizes, not as a coverage claim uniform across schemes.

\paragraph{What the envelope cannot do.} \IntervalsExclZero{} of \Nsch{}
intervals exclude zero: the tool can never state that a scheme \emph{will}
cost accuracy. And \BigLossBelow{} of \BigLoss{} observed losses worse
than 3pp fell \emph{below} their interval floor---when real damage occurs,
the interval usually misses it. This is the single most important caveat
in the paper. The obvious mitigation, an asymmetric band, was tried and
lost on both coverage and width where it was finite (\S\ref{sec:audit}),
so there is no drop-in fix on offer.

\section{The measured left tail}\label{sec:tail}

The corpus contains only configurations someone chose to publish. To measure
what it structurally cannot show, we rented an A100 and quantized real models
under deliberately-faulted recipes.

\paragraph{Result.} Worst measured loss \(\AdvWorst{}\)pp, against
\(\WorstWfour{}\)pp as the worst loss anywhere in the published corpus.
\AdvBadOver{} of \AdvBadRows{} faulted rows lost 3pp or more; separately,
\AdvCtrlOver{} of \AdvCtrlRows{} rows from our correctly-configured control arm
did too. The dominant single effect was wrong scale granularity---per-tensor
instead of per-group.

\paragraph{Scope, stated plainly.} This tail was measured only on models of
\AdvMaxB{}B and below. It establishes that catastrophic damage is reachable; it
is not a measured bound for larger models.

\paragraph{A confound we found and did not resolve.} Re-running MMLU under
the published protocol (5-shot, letter-scored) recovers accuracies in the
range published on the cards: our Qwen2.5-0.5B-Instruct scores
\BaseQwenSmall{} against the Instruct card's \PubQwenSmall{} (the base-model
card gives \PubQwenSmallBase{}), and our Qwen2.5-1.5B-Instruct scores
\BaseQwenMid{}. The harness question closed. Our control run then compared
Instruct-model checkpoints against Red Hat's \emph{base-model} W4A16
publications, because no Instruct W4A16 checkpoints are published: our
correctly-configured Instruct control is \CtrlRatio\(\times\) more damaging
than the base-model published W4A16 (mean gap \CtrlGap{}pp,
\(z = \CtrlZmin\) and \(\CtrlZmax\)). This is not like-for-like---Instruct
models could plausibly be more fragile under 4-bit quantization than their
base counterparts---so the ratio is indicative rather than a clean estimate.
A matched Instruct-vs-Instruct control is deferred future work. The pattern
we observed is consistent with approximately \CalibRatio\(\times\) less
calibration data than production defaults, though the attribution is not
independently tested. Consequently, excess-over-control measures
\emph{bad recipe versus mediocre recipe}, not versus production. It is an
indicative anchor, not a bound, and we decline to publish a corrected point
estimate: the correction moves by \ProtoSens{}pp depending on the anchor set,
and we have only \ExcessN{} rows from two models.

\section{Limitations}\label{sec:limits}

\paragraph{The guarantee is conditional, and absent where it matters most.}
Split conformal gives exact coverage under exchangeability between calibration
and test. Our calibration set is the set of configurations a publisher chose to
release---a selection event. Conditional on selection, calibration points are
not exchangeable with an arbitrary test point and the guarantee does not
transfer~\cite{barber2023beyond,jin2024focal}. It holds for a new checkpoint
drawn from the same publication process. It carries \textbf{no guarantee} for a recipe the
reader tuned themselves and that no publisher would have released---which is the case
a practitioner most wants covered.

\paragraph{Accuracy is not the whole of deployment risk.}
\citet{tong2026compression} show compression frequently decouples
accuracy from uncertainty: a model can hold accuracy while its prediction sets
inflate. A scheme this work calls low-risk on accuracy may still degrade
uncertainty quality. We measure none of that.

\paragraph{One publisher, one toolchain, and why.} See \S\ref{sec:data}. Our own
control arm (\S\ref{sec:tail}) is direct evidence that practitioner and
calibration volume can dominate scheme choice. We measured how confining this
is: surveying \PubChecked{} publishers and \PubCards{} quantized model cards,
\PubOthersPaired{} of \PubOthersCards{} cards from the other
\PubOthersCount{} publishers carry paired before/after evaluations, \PubIntelPaired{} of them Intel's, so others
do publish them. What the other publishers almost never provide is a
\emph{Recovery} column, which is what lets each row be checked against
\(100 \times \text{after}/\text{before}\) and rejected on disagreement: of
their rows, \PubOthersVerified{} pass that gate, all from
\PubOthersVerifiableCards{} card, against \PubCtrlVerified{} from the control
publisher. The corpus is single-publisher not because others publish nothing,
but because the control publisher's cards are, by a wide margin, the ones
published in a self-verifying form.

\paragraph{What can and cannot be re-run.} The supplement recomputes every tagged figure from the data files it ships. It cannot
re-run the faulted-configuration experiments (they need a GPU; the recorded results are included), the harvest from raw model
cards (the cards are not redistributed), or the prospective-card checks (those files are not included), and re-running the
low-rank comparison mentioned in the supplement needs a checkout of BenchPress's released code.

\paragraph{An attempt at re-running the calibration-data test.} We
attempted to reproduce Red Hat's published W4A16 checkpoints for two of
the corpus's most-measured base models, Qwen2.5-0.5B and Qwen2.5-1.5B,
using the \texttt{llm-compressor} recipe at its documented defaults, in
order to test whether the roughly \CalibRatio\(\times\) calibration-data
ratio between our own control and production defaults accounts for the
observed excess damage. Neither of these two cards pins a library version
or names a calibration dataset; both record only the resulting weight
representation (symmetric per-group INT4, group size 128 on 1.5B and
group size 64 on 0.5B). Neither attempt succeeded. A first run at the
current \texttt{llm-compressor} + \texttt{compressed-tensors} +
\texttt{transformers} stack produced quantized models scoring at or below
four-way random chance on 5-shot MMLU (\AnomQuantMinPct--\AnomQuantMaxPct\% absolute), well past the
worst outcome anywhere in the corpus or the control arm, which we treat
as a stack-level anomaly rather than a calibration-data finding. A second
run at an era-appropriate pin inferred from the card wording
(\texttt{llmcompressor==0.4.1}, \texttt{transformers==4.44.2}) could not
resolve to a working import: the pinned \texttt{llmcompressor} requires
a class that was only added to \texttt{transformers} in a later release,
so the pin does not exist as a working combination. The calibration-data
hypothesis therefore remains what the Results section already calls it,
consistent with the observation but not independently tested. That the
recipe recorded on these two cards is not sufficient to reproduce even
the working-stack question, let alone the answer, is directly
consistent with the sparsity of documentation reported in \S\ref{sec:data}.

\paragraph{Rows are not independent.} Up to \MaxBenchmarksPerRun{} benchmark rows come from a
single quantization run. We mitigate with group-conditional calibration and the
support floor, and note that a bootstrap over the \ProspClusN{} quantized checkpoints gives \([\ProspClusLo, \ProspClusHi]\)\%, close to the row-level interval,
but the unit of analysis remains a row. The leave-family-out coverage measurement has the
same dependence: each row is scored once under each of \PairsPerRow{} calibration families, so
the \PooledPairs{} evaluations are not independent, and every coverage judgment in this paper
resamples whole checkpoints rather than counting evaluations.

\section{Audit}\label{sec:audit}

We record errors found in our own work, including those that changed published
claims. This section exists because the alternative---reporting only what
survived---would misrepresent how the result was obtained.

\paragraph{Errors that changed a published number.}
\begin{itemize}
  \item \emph{The training corpus originally counted Llama-3.1, Llama-3.2
    and Llama-3.3 as three separate families.} Meta's own model cards
    state that Llama-3.2 1B/3B were pretraining-time distillations from
    Llama-3.1 8B/70B logits, that Llama-3.2 11B/90B are Llama-3.1
    plus a vision adapter, and that Llama-3.3-70B is a Llama-3.1-70B
    fine-tune. Under the mechanistic definition adopted in
    \S\ref{sec:data} (a distillation is a member of its teacher's
    family; a fine-tune inherits the base's family) they therefore
    belong to a single Llama-3 family, and the corpus has
    \Nfam{} training families rather than 8. Leave-one-family-out is
    now run over \Nfam{} folds rather than 8. The prospective
    coverage figures are unchanged (the strict subset never contained
    a Llama-3.2 or Llama-3.3 checkpoint), but every LOFO-derived
    figure moves: the aggregate LOFO coverage across held-out families
    changes from 87.3\% to 88.2\%, and the previously-worst held-out
    fold Llama-3.3 (85.1\%) folds into the merged Llama-3 fold at
    92.2\%. Per-family LOFO on the merged fold is reported so this
    weakness does not disappear from view.
  \item \emph{A retracted headline, then re-partitioned.} An early claim
    that all \ProspClusN{} held-out checkpoints came from families never
    used in construction was wrong. Under the definition adopted in
    \S\ref{sec:data}, \ProspInFamCkpts{} of the \ProspClusN{} are
    derivatives of a training family and \ProspOutFamCkpts{} are new
    pretraining runs; the point estimates on both sides land near nominal
    but the difference between them is not resolvable at this sample size
    (\S\ref{sec:results}).
  \item \emph{LOFO means different things for different schemes.}
    Under the shipped calibration path, when a leave-one-family-out
    calibration family has fewer than 9 rows of a scheme, the code
    keeps the training-set half-width for that scheme and labels the
    result as if it were calibrated. Under the \Nfam-family corpus
    this fallback fires on \LofoFallbackPairs{} of \LofoTotalPairs{}
    LOFO evaluations (\LofoFallbackRate\%), and the rate is highly
    scheme-specific: \LofoFallRateFp\% for fp8, \LofoFallRateNvfp\%
    for nvfp4, \LofoFallRateWsixteen\% for w8a16, \LofoFallRateFpdyn\%
    for fp8\_dynamic, \LofoFallRateWint\% for w8a8\_int, and
    \LofoFallRateWfour\% for w4a16 (never). A LOFO figure quoted at
    the pooled level is therefore an average over a construction that
    is genuinely held-out for w4a16, mostly held-out for fp8\_dynamic
    and w8a8\_int, and predominantly in-sample for fp8, nvfp4 and
    w8a16. On the held-out subset alone (fallback rows excluded),
    pooled two-sided coverage moves from \PooledCellCov\% to
    \LofoHeldoutCov\% and one-sided from \PooledOneSided\% to
    \LofoHeldoutOneSided\%. We disclose here rather than exclude:
    the shipped pipeline reports the pooled figure, and the reader
    should know what the pool contains.
  \item \emph{Independent implementations shared a float comparison.} The prospective-coverage figure sits on a floating point boundary at one row: gemma-3-1b-it W4A16 on TruthfulQA has an observed delta of exactly $+1.40$pp and a scheme upper bound of $+1.40$pp; in real arithmetic that is a closed-interval hit, but in float arithmetic the computed upper bound is $+1.3999999999999997$, so the comparison $\text{delta} \le \text{hi}$ returned False. All three reimplementations (the pipeline, the standard-library re-check, and the audit-side recomputation) agreed on the same 118/131 because they shared the strict comparison; the agreement was consistency, not correctness. The comparison sites were patched to add a $10^{-9}$ tolerance on each bound, and the strict figure is now \ProspInside/\ProspN. This is a real limit of independent reimplementation as a check: when two implementations share an assumption, they will agree on results the assumption is wrong about.
  \item \emph{GPQA was five protocols under one label.} The label
    \texttt{gpqa} in the harvested corpus was found to conflate at
    least five distinct GPQA evaluation protocols. Inspection of every
    RedHatAI card that reported a GPQA figure gives the following
    verbatim card labels and row counts, which we now carry under
    protocol-specific labels rather than one collapsed \texttt{gpqa}:
    \texttt{GPQA (0-shot)} $\to$ \texttt{gpqa\_main} (15 rows);
    \texttt{GPQA (Acc-Norm, 0-shot)} $\to$ \texttt{gpqa\_main\_norm}
    (8 rows); \texttt{GPQA CoT main (5-shot)} $\to$
    \texttt{gpqa\_main\_cot\_5shot} (2 rows);
    \texttt{GPQA (Diamond, 0-shot)} $\to$ \texttt{gpqa\_diamond}
    (8 rows); \texttt{GPQA CoT diamond (5-shot)} $\to$
    \texttt{gpqa\_diamond\_cot\_5shot} (1 row). The first-table-wins
    dedup that produced the earlier \texttt{gpqa} label kept whichever
    protocol each card listed first, so a single label mixed 0-shot
    unnormalised, 0-shot Acc-Norm, and 5-shot CoT scoring against both
    Main and Diamond subsets. We relabel rather than drop: the
    measurements are real under their respective protocols. Four rows
    lie under the modeling-set filter's chance floor
    (\texttt{acc\_before}$\geq$\MinAccBefore\%) and are excluded from the fitted
    set exactly as before; the split does not add or remove
    modelling-set rows. Every GPQA row now carries a card-verified
    protocol label; the harvester should adopt the same split so
    future corpora do not re-introduce the conflation.
  \item \emph{Nine pre-registered prospective repos were dropped by a
    name-gate.} The prospective set was pre-registered as
    \ProspRegisteredRepos{} repositories. The harvester requires a
    \texttt{<n>B} parameter-count token in the checkpoint name; the
    \ProspNameGatedRepos{} pre-registered repos whose names lack that
    token (phi-4 x3, Phi-4-mini x2, Mistral-Nemo x2, Devstral x2) are
    rejected by the name-gate and never enter the harvested corpus.
    The predictor is frozen and does not use size, so the gate is not
    load-bearing on the method; it is load-bearing on which of the
    pre-registered repos actually get scored. The paper reports on the
    \ProspN-row frozen-and-scored set (\ProspInside/\ProspN{} =
    \ProspCov\%). An external audit reports that, if the
    \ProspNameGatedRepos{} gated cards were harvested and scored under
    the same protocol, the with-them figure would be
    \ProspWithGatedInside/\ProspWithGatedTotal{} = \ProspWithGatedCov\%;
    that figure is not locally reproducible in this repository because
    the raw cards for those nine repos are not shipped. We report the
    frozen-set figure as our headline and disclose the with-them
    figure as an external check that the frozen-set number is not
    materially inflated by the drop.
  \item \emph{The family classifier is enumerated.} The in-family /
    out-of-family split of the \ProspClusN{} strict prospective checkpoints
    was decided by inspecting each card and matching against the
    training-family list; the code that emits the split is a regex specific
    to this set (\texttt{DeepSeek-R1-Distill-(Llama|Qwen)}). It is not a
    general implementation of the \S\ref{sec:data} definition. A
    prospective checkpoint outside this pattern (e.g. a distillation from
    another base, a fine-tune with an idiosyncratic name) would need a
    manual re-classification. The definition itself is fixed; the code
    path applies it to this set only.
  \item \emph{Two parser bugs.} A benchmark label typo cost 11 training rows;
    separately, \texttt{MATH-500} was matched as \texttt{Math-Lvl-5},
    contaminating 7 test rows with scores from a different benchmark.
  \item \emph{A centring mismatch.} Calibration measured residuals about the
    stratum mean while intervals centred on the scheme mean, mis-sizing every
    calibrated width. Found only by independent reimplementation.
  \item \emph{A column-order fault.} One card template orders columns
    \([\text{recovery}, \text{base}, \text{quant}]\), which our parser turned into a spurious
    \(-68\)pp delta. The integrity gate rejected 10 of 11 rows on that card,
    passing only the one where \(\text{base} = \text{quant}\) made the check
    degenerate.
  \item \emph{A retracted single-publisher claim.} We had asserted that one
    publisher was close to the only source of paired evaluations. A census of
    \PubChecked{} publishers found \PubOthersPaired{} paired-evaluation cards
    outside the control publisher, \PubIntelPaired{} of them Intel's; what
    differs is whose rows can be verified (\S\ref{sec:limits}).
  \item \emph{A bootstrap on the wrong statistic.} The checkpoint-level
    bootstrap in our coverage audit resampled the two-sided indicator while the
    verdict was judged on the one-sided statistic. An external commenter asked
    for it to wrap the one-sided figure, which exposed the mismatch
    (\S\ref{sec:method}).
  \item \emph{Inflated pair counts.} An earlier version of the tool reported
    the number of (test row, calibration family) evaluations as if it were the
    number of rows: with \Nfam{} families each row is scored \PairsPerRow{} times, so
    \PooledPairs{} evaluations came from only \PooledRows{} rows. Every
    row-count claim was recomputed on \texttt{distinct\_rows}.
  \item \emph{Two levels judged differently.} Cells were judged on the
    one-sided coverage statistic and refused when the point estimate fell
    below the threshold, but schemes had been judged on the two-sided
    statistic and refused only when the entire bootstrap interval fell below
    it. A single classifier now judges both levels the same way, and the
    envelope's refused list is a computed copy of the ranking tool's, checked
    by a test.
\end{itemize}

\paragraph{A regression we nearly shipped.} An audit indicated that a raw
asymmetric empirical band covered better than the shipped symmetric conformal
interval (\BandEmpCov\% versus \BandConfCov\%). Since quantization damage is left-skewed, a
symmetric band is obviously the wrong shape, and we were about to switch. On
measuring first, the empirical band returns an \emph{infinite} interval on \InfShare\%
of rows, because a two-sided empirical index requires more calibration points
than are often available. An infinite interval covers by construction. On rows
where both are finite, the original wins on coverage \emph{and} width. The
change was reversed.

\paragraph{An overclaimed guarantee.} The artifact was described as carrying a
coverage guarantee without qualification. Checking the method against the
conformal literature rather than against our own earlier work showed the
guarantee is conditional on exchangeability that selection breaks. Corrected in
\S\ref{sec:limits}.

\paragraph{Provenance discipline.} Two parser fixes were informed by inspecting
rows that were already in the test set. Rather than silently retain them, we
record which fixes were test-informed and exclude the affected family from the
strict headline.

\paragraph{An attempted extension of the adversarial arm to 3B.} On a fresh
40\,GB A100 with a fresh dependency install, we re-ran the adversarial
sweep to (i) verify that yesterday's numbers reproduce and (ii) add a
Qwen2.5-3B-Instruct row. The regression check passed: base evaluations
on Qwen2.5-0.5B-Instruct and SmolLM-135M-Instruct reproduced byte-for-byte,
deterministic-quantization recipes (per-tensor RTN, w8a8 without smoothing)
reproduced byte-for-byte, and stochastic GPTQ recipes stayed within
\(\leq 0.85\)pp of yesterday, consistent with fp16 Hessian-accumulation
noise. The 3B extension did not land: base evaluations on Qwen2.5-3B-Instruct
completed, but all five quantization recipes hit
\texttt{torch.cuda.OutOfMemoryError} during the paired-evaluation forward
pass at batch size 2, not during quantization itself; the fp16 model
weights (\(\approx 6\)\,GB), PyTorch's caching allocator fragmentation
from unloaded prior states, and activation memory for a batch of 2 long
sequences together exceeded the 40\,GB budget with roughly 10\,GB of
free memory needed for the next allocation. A single 80\,GB card would
be expected to accommodate this, so the constraint is hardware, not
methodology. The adversarial arm therefore remains bounded at
\(\leq\AdvMaxB\)B as before; \S\ref{sec:limits} states this scope
directly.

\paragraph{Mechanical verification.} The figures this paper reports as
results are macros generated from a claims registry that recomputes each value
from source artifacts, and the accompanying code fails a test if an
outward-facing document contains an untraced numeric literal. The exceptions
are historical figures in the error list above (row counts and the
\(-68\)pp delta), which are quoted from the provenance record rather than
recomputed.

\section{Conclusion}

Per-model prediction of quantization accuracy damage from card-level
features does not work on published data, and the reason is largely
measurement noise rather than model capacity; predictors that use
checkpoint-level statistics were not tested and remain open. What does
work is narrower: a per-scheme calibrated envelope that is validated at
\(\ProspCov\%\) two-sided coverage on unseen checkpoints for \ProspSchemes{}
of the \Nsch{} schemes. In-family and out-of-family point estimates both
land near nominal (\ProspInFamCov\% and \ProspOutFamCov\% respectively),
but the sample size does not distinguish them, and Gemma-3, the largest
out-of-family sub-group, covers at \ProspGemThreeCov\%, \ProspGemThreeGap{}pp below
nominal. The envelope also says which of its own cells it cannot judge. Each cell is trusted; or insufficient evidence, in which case the interval is still
shown but flagged and the scheme demoted; or refused, in which case the interval is withheld. Today \NCellsInsuff{} of
\NCellsTotal{} cells are insufficient evidence; on this corpus the refusal
state does not fire (\S\ref{sec:method}). The more useful contribution may be the limitation
itself---the published record of quantized models is outcome-truncated, and we
measured, on deliberately faulted recipes for models up to \AdvMaxB{}B, losses more than four times worse than anything it contains. A risk
estimate built on that record is a floor on risk, not a ceiling, and should say
so.

\bibliographystyle{plainnat}
\bibliography{refs}

\end{document}
